\chapter{周长与面积}\label{ch:g5:areas}

四年级时已经数过方格（见\cref{ch:g4:measure}）；本章
首次得到面积\emph{公式}：长乘宽，适用于
长方形，并学习面积单位 cm$^2$ 和 m$^2$。
三角形和其他图形的面积公式，将在\cref{ch:g6:measure}
及\cref{ch:g7:areas}中逐步学习。

\section{两种不同的度量}

\begin{example}[周长与面积]\label{ex:g5:areas:contrast}
园丁买篱笆要用到\emph{周长}，买草籽要用到
\emph{面积}。两者并不同步变化：把长方形拉得又细又长，
面积可以保持不变，而周长却会增大——比较
$6 \times 2$ 与 $4 \times 3$ 的两个长方形
（面积都是 $12$，周长分别是 $16$ 和 $14$）。
\end{example}

\section{长方形的面积公式}

\begin{proposition}[长方形的面积]\label{prop:g5:areas:rectangle}
设长方形的长为 $L$、宽为 $w$（两者单位相同），则其面积为
\[
A = L \times w ,
\]
面积以平方单位计：cm$^2$（边长 $1$ cm 的正方形），m$^2$（
边长 $1$ m 的正方形）、\dots
\end{proposition}

\begin{proof}[用计数说明理由]
用单位正方形铺满长方形：共 $w$ 行，每行 $L$ 个正方形，
共 $L \times w$ 个正方形——正是
\cref{ch:g3:mult}中的乘法长方形。
\end{proof}

\begin{omfigure}
\begin{tikzpicture}[scale=0.62]
  \draw[gray!50, thin] (0,0) grid (6,3);
  \draw[very thick, omDef] (0,0) rectangle (6,3);
  \fill[omThm!30] (0,2) rectangle (1,3);
  \draw[omThm, thick] (0,2) rectangle (1,3);
  \node[right, font=\small, omThm] at (1.15,2.5) {$1$ cm$^2$};
  \node[below, font=\small] at (3,-0.15) {$6$ cm};
  \node[left, font=\small] at (0,1.5) {$3$ cm};
  \node[font=\small] at (3.9,1.2) {$6 \times 3 = 18$ cm$^2$};
\end{tikzpicture}

{\small 三行，每行六个边长为 1 cm 的小方格：公式
$L \times w$ 就是对方格的计数，一次写定，永远适用。}
\end{omfigure}

\begin{example}\label{ex:g5:areas:rectangle}
一块地毯长 $2.5$ m、宽 $2$ m，面积为 $2.5 \times 2 = 5$ m$^2$。
一枚邮票长 $3$ cm、宽 $2.4$ cm，面积为 $3 \times 2.4 = 7.2$
cm$^2$。同一公式适用于任何单位，只要长与宽用的是
\emph{同一}种单位。
\end{example}

\begin{example}[平方单位按 100 换算]\label{ex:g5:areas:units}
$1$ m $= 100$ cm，但是 $1$ m$^2 = 100 \times 100 = 10\,000$ cm$^2$：
一平方米就是一个 $100$ 乘 $100$ 的平方厘米方格阵。面积
单位按百倍递增，而不是十倍。
\end{example}

\section{组合图形}

\begin{method}[切割、相加、相减]\label{met:g5:areas:composite}
对于由长方形组成的图形（如 L 形、T 形、框形）：
\begin{enumerate}
  \item 把图形切割成若干长方形，或补全为一个大长方形；
  \item 用公式计算每个长方形的面积；
  \item 把各部分面积相加，或者从
        大长方形中减去缺口部分的面积；
  \item 与方格的粗略计数相互核对。
\end{enumerate}
\end{method}

\begin{example}\label{ex:g5:areas:composite}
一间 L 形的房间：一个 $6$ m $\times$ $4$ m 的长方形
被截去一个 $2$ m $\times$ $2$ m 的角。
\[
\text{面积} = 6 \times 4 - 2 \times 2 = 24 - 4 = 20 \text{ m}^2 .
\]
也可以把 L 形切成一个 $6 \times 2$ 的长条和一个 $4 \times 2$ 的长条：
$12 + 8 = 20$ m$^2$——两条路，一个答案。
\end{example}

\begin{example}[长方形的一半]\label{ex:g5:areas:halfrect}
沿对角线把长方形剪开，得到两个面积相同的
三角形，每个恰好是长方形的\emph{一半}。直角边分别为
$6$ cm 和 $4$ cm 的直角三角形，面积因此为 $\frac{6 \times 4}{2} = 12$
cm$^2$——这是值得牢记的图示，
它将在\cref{ch:g6:measure}中变成公式。
\end{example}

\begin{omfigure}
\begin{tikzpicture}[scale=0.62]
  \draw[gray!50, thin] (0,0) grid (6,4);
  \draw[very thick] (0,0) rectangle (6,4);
  \fill[omDef!25] (0,0) -- (6,0) -- (0,4) -- cycle;
  \draw[very thick, omDef] (6,0) -- (0,4);
  \node[font=\small] at (1.7,1.2) {一半};
  \node[font=\small] at (4.3,2.8) {一半};
\end{tikzpicture}

{\small 对角线把 $6 \times 4$ 的长方形分成两个相等的
三角形，每个含 $12$ 个小方格。}
\end{omfigure}

\section{练习}

\begin{exercise}[$\star$]\label{exo:g5:areas:1}
计算长方形的面积和周长，其尺寸为 $8$ cm $\times$
$5$ cm。哪个结果的单位是 cm，哪个是 cm$^2$？
\end{exercise}

\begin{exercise}[$\star$]\label{exo:g5:areas:2}
求下列图形的面积：边长为 $9$ cm 的正方形；长方形
$12$ m $\times$ $7$ m；长方形 $4.5$ cm $\times$ $6$ cm。
\end{exercise}

\begin{exercise}[$\star$]\label{exo:g5:areas:3}
一个长方形的面积为 $63$ cm$^2$，长为 $9$ cm。求它的宽，
再求它的周长。
\end{exercise}

\begin{exercise}[$\star$]\label{exo:g5:areas:4}
在方格纸上画两个面积都为 $24$ 个小方格的不同长方形，
并分别计算它们的周长。面积相同 --- 周长也相同吗？
\end{exercise}

\begin{exercise}[$\star$]\label{exo:g5:areas:5}
换算下列单位：$3$ m$^2$ 化为 cm$^2$；\quad $50\,000$ cm$^2$ 化为 m$^2$。
（回顾 \cref{ex:g5:areas:units}：按百进行换算！）
\end{exercise}

\begin{exercise}[$\star$]\label{exo:g5:areas:6}
一个 T 形图形由一个 $8 \times 2$ 的水平长条
放在一个 $2 \times 5$ 的竖直长条之上（单位：cm）。求其面积，可将它看作两个
长方形，分别计算后相加。
\end{exercise}

\begin{exercise}[$\star$]\label{exo:g5:areas:7}
一个画框由一个 $30$ cm $\times$ $20$ cm 的长方形
挖去一个 $24$ cm $\times$ $14$ cm 的长方形窗口制成。
求画框所用木材的面积是多少？
\end{exercise}

\begin{exercise}[$\star$]\label{exo:g5:areas:8}
求两条直角边分别为 $8$ cm 和 $6$ cm 的直角三角形的面积
（它恰为一个长方形的一半，见\cref{ex:g5:areas:halfrect}）。
\end{exercise}

\begin{exercise}[$\star$]\label{exo:g5:areas:9}
一块长方形菜地长 $7$ m、宽 $4$ m。种子每平方米
$2$ 元，篱笆每米 $3$ 元。求种子的费用，
再求篱笆的费用。
\end{exercise}

\begin{exercise}[$\star\star$]\label{exo:g5:areas:10}
一条走廊的地面长 $12$ m、宽 $2$ m，需要用
边长为 $50$ cm 的正方形瓷砖铺满。需要多少块瓷砖？
（先换算单位，或沿每个方向数出瓷砖的块数。）
\end{exercise}

\begin{exercise}[$\star\star$]\label{exo:g5:areas:11}
将一个 $4$ cm $\times$ $3$ cm 的长方形的各边扩大为两倍，周长会怎样变化？
面积呢？（扩大前后的周长和面积都要计算 ---
两个结果不同！）
\end{exercise}
